\chapter{Chapter05 — Sparse Polynomial Approximation}
%%% generated by the blueprint pipeline from FABL.Chapter05.SparsePolynomialApproximation
\section{Overview}

Declarations of the Lean module \texttt{FABL.Chapter05.SparsePolynomialApproximation}.

\section{Declarations}

\begin{theorem}[Positivity of the Fourier one-norm]
\lean{FABL.fourierOneNorm_pos_of_ne_zero}
\label{FABL.fourierOneNorm_pos_of_ne_zero}
\leanok
\uses{FABL.abs_apply_le_fourierOneNorm, FABL.fourierOneNorm}
\difficulty{3}
Let $n$ be a nonnegative integer and let $f$ be a real-valued function on the sign cube
$\\{-1,1\\}^n$. If $f$ is not identically zero, then its Fourier one-norm is strictly
positive.
\end{theorem}

\begin{theorem}[Empirical mean concentration for finite laws]
\lean{FABL.measure_finitePMFEmpiricalMean_sub_integral_ge_le}
\label{FABL.measure_finitePMFEmpiricalMean_sub_integral_ge_le}
\leanok
\uses{FABL.finiteUniformEmpiricalMean, FABL.independentProductPMF, FABL.independentProductPMF_toMeasure}
\difficulty{5}
Let $\\Omega$ be a finite nonempty set equipped with a measurable space in which every
singleton is measurable. Let $p$ be a probability mass function on $\\Omega$, and let
$g:\\Omega\\to\\mathbb R$ take values in $[-1,1]$. For a positive integer $m$ and
$\\varepsilon\\ge 0$, draw $m$ independent samples from $p$. The probability that their
empirical mean differs from the integral of $g$ with respect to $p$ by at least
$\\varepsilon$ is at most $2\\exp(-m\\varepsilon^2/2)$.
\end{theorem}

\begin{definition}[Fourier atom probability distribution]
\lean{FABL.fourierAtomPMF}
\statementsource{lmn-constant-depth-fourier}{
  pdf-63759bd909d5-p004-b003,
  pdf-63759bd909d5-p004-b004
}
\statementsource{Analysis-of-Boolean-Functions-by-Ryan-ODonnell}{
  pdf-ffbe9c7b3306-p027-b007,
  pdf-ffbe9c7b3306-p027-b008,
  pdf-ffbe9c7b3306-p117-b005
}
\label{FABL.fourierAtomPMF}
\leanok
\uses{FABL.fourierCoeff, FABL.fourierOneNorm, FABL.fourierOneNorm_pos_of_ne_zero}
For a nonzero function $f:\{-1,1\}^n\to\mathbb R$, define a probability distribution on
subsets $S\subseteq\{1,\ldots,n\}$ by assigning mass $|\widehat f(S)|/\|\widehat f\|_1$
to $S$, where $\widehat f(S)$ is the Fourier coefficient and $\|\widehat
f\|_1=\sum_T|\widehat f(T)|$.
\end{definition}

\begin{definition}[Signed Fourier atom]
\lean{FABL.signedFourierAtom}
\statementsource{lmn-constant-depth-fourier}{pdf-63759bd909d5-p004-b002}
\statementsource{Analysis-of-Boolean-Functions-by-Ryan-ODonnell}{pdf-ffbe9c7b3306-p117-b005}
\label{FABL.signedFourierAtom}
\leanok
\uses{FABL.fourierCoeff, FABL.monomial}
For a function $f:\{-1,1\}^n\to\mathbb R$, a subset $S\subseteq\{1,\ldots,n\}$, and an
input $x\in\{-1,1\}^n$, define the signed Fourier atom as the sign of $\widehat f(S)$
multiplied by the monomial $\prod_{i\in S}x_i$.
\end{definition}

\begin{theorem}[Signed Fourier atom bounds]
\lean{FABL.signedFourierAtom_mem_Icc}
\proofsource{Analysis-of-Boolean-Functions-by-Ryan-ODonnell}{pdf-ffbe9c7b3306-p117-b005}
\label{FABL.signedFourierAtom_mem_Icc}
\leanok
\uses{FABL.fourierCoeff, FABL.monomial_sq, FABL.signedFourierAtom}
\difficulty{4}
Let $n\ge 0$, let $f:\{-1,1\}^n\to\mathbb R$, let $S\subseteq\{0,\ldots,n-1\}$, and let
$x\in\{-1,1\}^n$. The product of the sign of the Fourier coefficient of $f$ at $S$ and
the monomial indexed by $S$, evaluated at $x$, lies in $[-1,1]$.
\end{theorem}

\begin{theorem}[Mean of a signed Fourier atom]
\lean{FABL.integral_signedFourierAtom_fourierAtomPMF}
\proofsource{Analysis-of-Boolean-Functions-by-Ryan-ODonnell}{pdf-ffbe9c7b3306-p117-b005}
\label{FABL.integral_signedFourierAtom_fourierAtomPMF}
\leanok
\uses{FABL.fourierAtomPMF, FABL.fourierCoeff, FABL.fourierOneNorm, FABL.fourierOneNorm_pos_of_ne_zero, FABL.fourier_expansion, FABL.monomial, FABL.signedFourierAtom}
\difficulty{4}
Let $n$ be a natural number, let $f:\{-1,1\}^n\to\mathbb R$ be nonzero, and fix
$x\in\{-1,1\}^n$. Choose a subset $S\subseteq\{1,\ldots,n\}$ with probability
proportional to the absolute value of the Fourier coefficient of $f$ at $S$, and
evaluate the corresponding Fourier monomial at $x$, multiplied by the sign of that
coefficient. The expected value is $f(x)$ divided by the sum of the absolute values of
all Fourier coefficients of $f$.
\end{theorem}

\begin{definition}[Sampled Fourier family]
\lean{FABL.sampledFourierFamily}
\statementsource{lmn-constant-depth-fourier}{pdf-63759bd909d5-p010-b002}
\statementsource{Analysis-of-Boolean-Functions-by-Ryan-ODonnell}{pdf-ffbe9c7b3306-p117-b005}
\label{FABL.sampledFourierFamily}
\leanok
Given $s\in\mathbb N$ and samples $S_i\subseteq\{1,\ldots,n\}$ for $i=1,\ldots,s$,
define the sampled Fourier family to be the set of subsets that occur among the samples.
\end{definition}

\begin{definition}[Sampled Fourier coefficient]
\lean{FABL.sampledFourierCoefficient}
\statementsource{lmn-constant-depth-fourier}{pdf-63759bd909d5-p010-b002}
\statementsource{Analysis-of-Boolean-Functions-by-Ryan-ODonnell}{pdf-ffbe9c7b3306-p117-b005}
\label{FABL.sampledFourierCoefficient}
\leanok
\uses{FABL.fourierCoeff, FABL.fourierOneNorm}
Given $s\in\mathbb N$, a function $f:\{-1,1\}^n\to\mathbb R$, samples
$S_i\subseteq\{1,\ldots,n\}$ for $i=1,\ldots,s$, and a subset
$T\subseteq\{1,\ldots,n\}$, define the sampled coefficient at $T$ to be $\|\widehat
f\|_1/s$ times the sum of $\operatorname{sgn}(\widehat f(T))$ over sample indices $i$
for which $S_i=T$.
\end{definition}

\begin{definition}[Sampled Fourier polynomial]
\lean{FABL.sampledFourierPolynomial}
\statementsource{lmn-constant-depth-fourier}{
  pdf-63759bd909d5-p010-b003,
  pdf-63759bd909d5-p010-b009
}
\statementsource{Analysis-of-Boolean-Functions-by-Ryan-ODonnell}{pdf-ffbe9c7b3306-p117-b005}
\label{FABL.sampledFourierPolynomial}
\leanok
\uses{FABL.sampledFourierCoefficient, FABL.sampledFourierFamily, FABL.sparseFourierApproximation}
Given $s\in\mathbb N$, a function $f:\{-1,1\}^n\to\mathbb R$, and samples
$S_i\subseteq\{1,\ldots,n\}$ for $i=1,\ldots,s$, define the sampled Fourier polynomial
to be the sparse Fourier approximation associated with the family of sampled subsets and
their sampled coefficients.
\end{definition}

\begin{theorem}[Sampled fourier polynomial apply]
\lean{FABL.sampledFourierPolynomial_apply}
\statementsource{lmn-constant-depth-fourier}{pdf-63759bd909d5-p010-b003}
\label{FABL.sampledFourierPolynomial_apply}
\leanok
\uses{FABL.finiteUniformEmpiricalMean, FABL.fourierCoeff, FABL.fourierOneNorm, FABL.monomial, FABL.sampledFourierFamily, FABL.sampledFourierPolynomial, FABL.signedFourierAtom}
\difficulty{5}
% dataset-pipeline: stub -- prose pending build_informal_statements.py
Statement pending regeneration.
\end{theorem}

\begin{theorem}[Sparsity of a sampled Fourier polynomial]
\lean{FABL.polynomialSparsity_sampledFourierPolynomial_le}
\label{FABL.polynomialSparsity_sampledFourierPolynomial_le}
\leanok
\uses{FABL.fourierCoeff_sparseFourierApproximation, FABL.fourierSupport, FABL.mem_fourierSupport, FABL.polynomialSparsity, FABL.sampledFourierFamily, FABL.sampledFourierPolynomial}
\difficulty{3}
Let $n,s$ be natural numbers, let $f$ be a real-valued function on the $n$-dimensional
sign cube, and choose a set of coordinates $S_i$ for each $imod s$. The number of
nonzero Fourier coefficients of the Fourier polynomial formed from the distinct sampled
sets $S_i$ is at most $s$.
\end{theorem}

\begin{theorem}[Strict exponential bound on the sign cube]
\lean{FABL.card_signCube_mul_hoeffding_bound_lt_one}
\proofsource{Analysis-of-Boolean-Functions-by-Ryan-ODonnell}{pdf-ffbe9c7b3306-p117-b005}
\label{FABL.card_signCube_mul_hoeffding_bound_lt_one}
\leanok
\uses{FABL.Sign}
\difficulty{5}
Let $n$ and $s$ be natural numbers with $n>0$, and let $L,\delta\in\mathbb R$ satisfy
$L>0$, $\delta>0$, and $4nL^2/\delta^2\le s$. The number of sign vectors in
$\{-1,1\}^n$, multiplied by $2\exp\!\left(-s(\delta/L)^2/2\right)$, is strictly less
than $1$.
\end{theorem}

\begin{theorem}[Exists fourier atom samples uniform approximation]
\lean{FABL.exists_fourierAtomSamples_uniformApproximation}
\statementsource{lmn-constant-depth-fourier}{
  pdf-63759bd909d5-p010-b005,
  pdf-63759bd909d5-p010-b006
}
\proofsource{Analysis-of-Boolean-Functions-by-Ryan-ODonnell}{pdf-ffbe9c7b3306-p117-b005}
\label{FABL.exists_fourierAtomSamples_uniformApproximation}
\leanok
\uses{FABL.card_signCube_mul_hoeffding_bound_lt_one, FABL.finiteUniformEmpiricalMean, FABL.fourierAtomPMF, FABL.fourierCoeff, FABL.fourierOneNorm, FABL.fourierOneNorm_pos_of_ne_zero, FABL.independentProductPMF, FABL.independentProductPMF_apply, FABL.integral_signedFourierAtom_fourierAtomPMF, FABL.measure_finitePMFEmpiricalMean_sub_integral_ge_le, FABL.sampledFourierPolynomial, FABL.sampledFourierPolynomial_apply, FABL.signedFourierAtom, FABL.signedFourierAtom_mem_Icc}
\difficulty{6}
% dataset-pipeline: stub -- prose pending build_informal_statements.py
Statement pending regeneration.
\end{theorem}

\begin{theorem}[Exists sparse polynomial uniform approximation of ne zero]
\lean{FABL.exists_sparsePolynomial_uniformApproximation_of_ne_zero}
\statementsource{lmn-constant-depth-fourier}{
  pdf-63759bd909d5-p011-b004,
  pdf-63759bd909d5-p011-b005
}
\proofsource{Analysis-of-Boolean-Functions-by-Ryan-ODonnell}{
  pdf-ffbe9c7b3306-p117-b004,
  pdf-ffbe9c7b3306-p117-b005
}
\label{FABL.exists_sparsePolynomial_uniformApproximation_of_ne_zero}
\leanok
\uses{FABL.exists_fourierAtomSamples_uniformApproximation, FABL.fourierOneNorm, FABL.polynomialSparsity, FABL.polynomialSparsity_sampledFourierPolynomial_le, FABL.sampledFourierPolynomial}
\difficulty{1}
% dataset-pipeline: stub -- prose pending build_informal_statements.py
Statement pending regeneration.
\end{theorem}

\begin{theorem}[Sparse uniform approximation]
\lean{FABL.exists_sparsePolynomial_uniformApproximation}
\proofsource{Analysis-of-Boolean-Functions-by-Ryan-ODonnell}{
  pdf-ffbe9c7b3306-p117-b004,
  pdf-ffbe9c7b3306-p117-b005
}
\label{FABL.exists_sparsePolynomial_uniformApproximation}
\leanok
\uses{FABL.exists_sparsePolynomial_uniformApproximation_of_ne_zero, FABL.fourierCoeff, FABL.fourierOneNorm, FABL.fourierSupport, FABL.polynomialSparsity}
\difficulty{4}
Let $n>0$, let $f$ be a real-valued function on the $n$-dimensional sign cube, and let
$\delta>0$. If $4n\|\widehat f\|_1^2/\delta^2\le s$, where $s$ is a nonnegative integer
and $\|\widehat f\|_1$ is the sum of the absolute values of the Fourier coefficients of
$f$, then there is a function $q$ on the cube with at most $s$ nonzero Fourier
coefficients such that $|f(x)-q(x)|<\delta$ at every point $x$.
\end{theorem}

\begin{theorem}[Nonvanishing of the real-valued function]
\lean{FABL.booleanFunction_toReal_ne_zero}
\label{FABL.booleanFunction_toReal_ne_zero}
\leanok
\uses{FABL.BooleanFunction, FABL.BooleanFunction.toReal, FABL.Sign, FABL.signValue_eq_neg_one_or_one}
\difficulty{4}
Let $n$ be a nonnegative integer and let $f$ be a sign-valued Boolean function on the
$n$-dimensional sign cube. The real-valued function obtained by viewing each output sign
as a real number is not the zero function; equivalently, its value is nonzero at every
input.
\end{theorem}

\begin{theorem}[Sparse polynomial threshold representation]
\lean{FABL.exists_polynomialThresholdRepresentation_sparsity_le_ceil}
\proofsource{Analysis-of-Boolean-Functions-by-Ryan-ODonnell}{
  pdf-ffbe9c7b3306-p117-b004,
  pdf-ffbe9c7b3306-p117-b005,
  pdf-ffbe9c7b3306-p117-b006,
  pdf-ffbe9c7b3306-p117-b007
}
\label{FABL.exists_polynomialThresholdRepresentation_sparsity_le_ceil}
\leanok
\uses{FABL.BooleanFunction, FABL.BooleanFunction.toReal, FABL.IsPolynomialThresholdRepresentation, FABL.exists_sparsePolynomial_uniformApproximation, FABL.fourierOneNorm, FABL.polynomialSparsity, FABL.signValue_neg_one, FABL.signValue_one, FABL.thresholdSign_of_neg, FABL.thresholdSign_of_nonneg}
\difficulty{4}
Let $n>0$ and let $f:\\{−1,1\\}^n\\to\\{−1,1\\}$. There is a real-valued function $p$ on
the cube whose threshold signs equal $f$ at every input, and whose Fourier support has
size at most
$\\left\\lceil4n\\left(\\sum_{S\\subseteq[n]}|\\widehat{f}(S)|\\right)^2\\right\\rceil$.
\end{theorem}

\begin{theorem}[Exists parity majority representation]
\lean{FABL.exists_parityMajorityRepresentation}
\statementsource{lmn-constant-depth-fourier}{
  pdf-63759bd909d5-p010-b003,
  pdf-63759bd909d5-p010-b004
}
\statementsource{Analysis-of-Boolean-Functions-by-Ryan-ODonnell}{pdf-ffbe9c7b3306-p117-b006}
\label{FABL.exists_parityMajorityRepresentation}
\leanok
\uses{FABL.BooleanFunction, FABL.BooleanFunction.toReal, FABL.booleanFunction_toReal_ne_zero, FABL.exists_fourierAtomSamples_uniformApproximation, FABL.finiteUniformEmpiricalMean, FABL.fourierCoeff, FABL.fourierOneNorm, FABL.fourierOneNorm_pos_of_ne_zero, FABL.monomial, FABL.sampledFourierPolynomial, FABL.sampledFourierPolynomial_apply, FABL.signValue_neg_one, FABL.signValue_one, FABL.signedFourierAtom, FABL.thresholdSign, FABL.thresholdSign_of_neg, FABL.thresholdSign_of_nonneg}
\difficulty{5}
% dataset-pipeline: stub -- prose pending build_informal_statements.py
Statement pending regeneration.
\end{theorem}
